\chapter{Testing and Results}

\section{Test strategy}

System validation relies on two levels of testing: unit tests verifying the isolated behavior of each business service, and integration tests verifying the end-to-end behavior of the main user journeys (creating a work order, raising an alert, the conversational assistant's response). The test suite covers 99 tests in total (46 unit and core tests, 53 backend integration tests), all passing against the live Docker stack.

\section{Machine Learning model evaluation}

\subsection{Validation methodology}

Chapter~3 already ruled out a random row split for this data, and the testing phase is where that design choice gets to prove itself. A model graded on a random 80/20 split would have been evaluated on machines it had already partly watched degrade, since both halves of a random split draw from the same fleet; any resulting score would measure memorisation of that fleet's specific trajectories rather than genuine ability to generalise to a machine never encountered before.

P3 is the model actually held to that standard in this chapter: machines 70--79 were withheld entirely before training began, so every figure reported for it below reflects performance on machines it never saw. It is, as a result, the only model here whose numbers can be read at face value; the rest stay gated on the label-quality issue covered next, independently of how their data was split.

\subsection{Label quality and validation status}
\label{sec:validation}

The risk was flagged early, in Chapter~1: the five failure sub-labels were never meant to double as predictors of the very flag they are computed from. What this testing phase did was turn that early caution into a confirmed finding, one Chapter~3 already treats as the platform's most significant data-quality issue: P1, P2, and P5 all read \texttt{tool\_wear} as an input while their own training labels are computed as a deterministic function of that same column, textbook target leakage. No amount of extra data, feature engineering, or algorithm swapping fixes a problem that lives in the label source rather than the model, so none was attempted.

What that confirmed leak does to the reported numbers is worth stating plainly: the $\sim$99\% AUC-ROC these three models post under group-holdout measures how faithfully each one reproduces an arithmetic identity it was never supposed to see stated so directly, not how well it predicts a real failure. Publishing that figure as a performance metric would hand stakeholders false confidence, which is why the model registry marks every model's \texttt{validation\_status} explicitly, and why the dashboard hides accuracy numbers for anything tagged \texttt{"leaked"} instead of showing a misleadingly high score.

Table~\ref{tab:validation-status} summarizes the validation status and production gate for each model.

\begin{table}[h]
\centering
\begin{tabularx}{\textwidth}{|X|X|X|X|}
\toprule
\textbf{Model} & \textbf{Algorithm} & \textbf{Production gate} & \textbf{Status} \\
\midrule
P1: Failure Probability   & XGBClassifier         & AUC-ROC $\geq 0.85$            & Leaked \\
\midrule
P2: Failure Type          & MultiOutputClassifier & Macro-F1 $\geq 0.70$           & Leaked (partial) \\
\midrule
P3: RUL Estimate          & XGBRegressor          & $R^2 \geq 0.40$, C-index $\geq 0.62$ & Group-holdout validated \\
\midrule
P4: Anomaly Detection     & Ensemble (IF+Z+AE)    & Precision on WO outcomes       & Unverified \\
\midrule
P5: Work Order Priority   & XGBClassifier         & F1 $\geq 0.70$                 & Leaked \\
\midrule
P6: Maintenance Schedule  & XGBClassifier         & F1 $\geq 0.70$                 & Unverified \\
\midrule
P7: Parts Demand          & Survival + Croston    & vs.\ naive baseline            & Unverified \\
\bottomrule
\end{tabularx}
\caption{Validation status and production gate per model}
\label{tab:validation-status}
\end{table}

\subsection{Validated model metrics}

For models with a reliable evaluation, Table~\ref{tab:metriques-ml} reports the measured metrics against the production gate.

\begin{table}[h]
\centering
\begin{tabularx}{\textwidth}{|X|X|X|X|}
\toprule
\textbf{Model} & \textbf{Metric} & \textbf{Result} & \textbf{Gate} \\
\midrule
P3: Remaining Useful Life & Mean Absolute Error (MAE)    & 14.95 cycles & $\leq 16.25 \times 1.15$ \\
\midrule
P3: Remaining Useful Life & Concordance index (C-index)  & 0.620        & $\geq 0.62$ \\
\midrule
P4: Anomaly Detection     & ROC-AUC (baseline)           & 0.83         & $>$ naive baseline \\
\midrule
P3: Prototype (Cox PH)    & C-index (research prototype) & 0.712        & reference only \\
\midrule
P3: Prototype (Weibull AFT) & C-index (research prototype) & 0.714      & reference only \\
\bottomrule
\end{tabularx}
\caption{Summary of model evaluation metrics (validated models only)}
\label{tab:metriques-ml}
\end{table}

The Cox and Weibull results in the last two rows come from a research prototype that applies censoring-aware survival modeling to the same held-out machine set as P3, and demonstrate that explicit handling of right-censored machines (those still running at the observation cutoff) produces a genuine improvement over the naive regression baseline (C-index 0.644). These results support the roadmap decision to migrate P3 toward a survival-model architecture once real run-to-failure data is available.

\subsection{Retraining strategy}

A production gate is enforced on every retrain: a new model artifact is only written to disk if it meets or exceeds the minimum threshold for its metric. This prevents a bad retrain from silently overwriting a good model.

The trigger strategy varies by model and depends on its validation status:

\begin{itemize}
\item \textbf{P1 and P5} are frozen until real technician-confirmed failure records replace the formula-derived labels. Retraining on more synthetic data would only reinforce the tautology.
\item \textbf{P2} is retrained on a manual trigger when a meaningful batch of new technician-confirmed failure-type labels accumulates, using a repeated group-holdout cross-validation pass to report mean and standard deviation of macro-F1 across folds.
\item \textbf{P3} is the only model scheduled for a regular cadence: monthly batch retraining, or triggered whenever a full machine lifecycle newly closes out and becomes available as a training sample. It is first in line for new regressor accuracy gating because it is the most trustworthy model in the pipeline.
\item \textbf{P4, P6, and P7} are retrained on an event-driven basis when sufficient real outcome data (confirmed anomaly flags, intervention outcomes, actual parts consumption) has accumulated to constitute a meaningful update signal.
\end{itemize}

\subsection{What real data would unlock}

The current platform is built on a solid architecture but is limited by the quality of available training data. Three concrete investments would change the situation significantly.

First, \textbf{technician-confirmed failure records}: when a technician closes a work order and records the root cause (TWF, HDF, PWF, OSF, or RNF), that observation is a label P1 and P2 can genuinely learn from. Accumulating 1\,000 or more confirmed events over six to twelve months is the minimum threshold for meaningful retraining.

Second, \textbf{real run-to-failure sequences for P3}: tracking fifteen to twenty machines continuously from healthy state to confirmed failure (or documented decommission) would provide P3 with real censored sequences, enabling a migration to Cox Proportional Hazards or Weibull AFT as the production model.

Third, \textbf{vibration and motor current sensors}: the five sensor inputs currently available cannot capture bearing degradation, the most common industrial failure mode. Vibration \textbf{Root Mean Square (RMS)} and kurtosis, together with motor current draw, are the signals that would make P1, P3, and P4 substantively more powerful, because they provide information that is genuinely independent of tool wear.

\subsection{Summary}

Of the seven models, only P3 (RUL) has evaluation figures that can be trusted: it is the one model trained and tested with proper group-holdout validation, achieving an MAE of 14.95 cycles and a C-index of 0.620 against its production gate. P1, P2, and P5 are formally flagged \texttt{"leaked"} because their labels are deterministic functions of an input feature, and the retraining strategy for every model is explicitly gated on accumulating real, technician-confirmed outcomes rather than more synthetic data.

\section{Use case demonstrations}

Three scenarios illustrate the end-to-end operation of the platform:
\begin{itemize}
\item \textbf{Creating and tracking an intervention order}: a technician receives an intervention order assigned by the Head Technician, views it from their space, and then closes it once completed, which updates the history of the machine concerned;
\item \textbf{Receiving a stock shortage alert}: when the P7 module detects that a part needed for a planned intervention will soon be unavailable, an alert is automatically generated and presented to the Head Technician along with an order proposal for validation;
\item \textbf{Querying the conversational assistant}: a technician asks the assistant about the troubleshooting procedure for a specific machine; the response combines the technical documentation with the machine's real-time health status.
\end{itemize}

\begin{figure}[H]
\centering
\includegraphics[width=\textwidth]{figures/ss_uc_intervention_1_request.png}
\caption{Step 1: technician submits the intervention request}
\label{fig:ss-uc-intervention-1}
\end{figure}

\begin{figure}[H]
\centering
\includegraphics[width=\textwidth]{figures/ss_uc_intervention_2_approval.png}
\caption{Step 2: Head Technician approves the request}
\label{fig:ss-uc-intervention-2}
\end{figure}

\begin{figure}[H]
\centering
\includegraphics[width=\textwidth]{figures/ss_uc_intervention_3_closure.png}
\caption{Step 3: technician closes the work order}
\label{fig:ss-uc-intervention-3}
\end{figure}

\begin{figure}[H]
\centering
\includegraphics[width=\textwidth]{figures/ss_uc_shortage.png}
\caption{End-to-end: stock shortage alert}
\label{fig:ss-uc-shortage}
\end{figure}

\subsection{Summary}

The intervention lifecycle (request, approval, closure), the stock-shortage-alert flow, and a conversational-assistant query were walked through end-to-end on the live platform: a request submitted by a technician reaches the Head Technician's queue with a linked work order, a projected shortage surfaces as an alert with a linked draft order, and a technician's question is answered using both documentation and the machine's real-time health, exactly as specified in Chapter~2.

\section{Difficulties encountered during testing}

The testing phase highlighted the importance of training data representativeness: the models were validated on a historical dataset simulating several dozen complete maintenance cycles, in order to cover a sufficient diversity of failure scenarios. The computation time of certain models (in particular the hybrid remaining-useful-life model) required optimizing model loading, which is cached in memory to avoid reloading from disk on every request.

A subtler difficulty involved the interaction between synthetic telemetry and real-time inference. When the platform was seeded with generated readings flagged \texttt{is\_synthetic = true}, every machine displayed identical sensor values and ML outputs, because the ML service was silently substituting fixed placeholder values when no real readings passed the synthetic filter. The fix replaced silent placeholder defaults with an explicit \texttt{telemetry\_available: false} flag, stopping the model call entirely when no real data exists and making the absence of data visible to the user rather than hidden behind a plausible-looking number.
